% =====================================================================
%  Bab I — Pendahuluan (rekonstruksi dari naskah PDF).
%  Diperiksa ulang oleh penulis dianjurkan (transkripsi dari PDF).
% =====================================================================

\section{Pendahuluan}

\subsection{Latar Belakang}
Perkembangan teknologi Internet of Things (IoT) telah mentransformasi
paradigma kehidupan rumah tangga modern secara fundamental. Atzori
\textit{et al.}~\cite{b5} mendefinisikan IoT sebagai sebuah paradigma di mana
objek-objek fisik dilengkapi dengan kemampuan komputasi, konektivitas
jaringan, dan kapasitas pertukaran data, sehingga memungkinkan interaksi dan
kolaborasi antarobjek tanpa intervensi manusia secara langsung. Menurut
laporan International Telecommunication Union (ITU)~\cite{b7}, jumlah perangkat
IoT yang terhubung secara global telah melampaui 15 miliar unit pada tahun
2024, dengan proyeksi mencapai 25 miliar unit pada tahun 2030. Pertumbuhan ini
didorong oleh meningkatnya adopsi ekosistem smart home yang mencakup perangkat
seperti smart TV, kamera pengintai berbasis IP (CCTV), smart speaker, smart
plug, smart lock, dan sensor lingkungan yang kesemuanya terhubung melalui
jaringan Wi-Fi domestik.

Al-Fuqaha \textit{et al.}~\cite{b6} dalam survei komprehensifnya
mengidentifikasi bahwa ekosistem IoT bergantung pada beragam protokol
komunikasi seperti MQTT (Message Queuing Telemetry Transport), CoAP
(Constrained Application Protocol), Zigbee, Z-Wave, dan Bluetooth Low Energy
(BLE), yang masing-masing memiliki karakteristik keamanan yang berbeda-beda.
Keragaman protokol ini, meskipun memberikan fleksibilitas dalam implementasi,
menciptakan kompleksitas tersendiri dalam pengelolaan keamanan jaringan rumah
tangga.

Seiring dengan meningkatnya adopsi perangkat IoT di lingkungan domestik,
ancaman keamanan siber terhadap jaringan rumah tangga turut mengalami eskalasi
yang signifikan. Salah satu insiden keamanan IoT paling menonjol dalam sejarah
adalah serangan Mirai Botnet pada tahun 2016. Antonakakis
\textit{et al.}~\cite{b1} mengungkapkan bahwa Mirai berhasil menginfeksi lebih
dari 600.000 perangkat IoT --- termasuk kamera IP, router, dan Digital Video
Recorder (DVR) --- dengan mengeksploitasi default credentials yang tidak diubah
oleh penggunanya. Perangkat-perangkat yang terinfeksi tersebut kemudian
digunakan untuk melancarkan serangan Distributed Denial of Service (DDoS) masif
yang melumpuhkan layanan internet berskala besar seperti Dyn DNS, Twitter,
Netflix, dan Reddit. Kolias \textit{et al.}~\cite{b10} menegaskan bahwa insiden
Mirai menjadi titik balik yang mengekspos kerentanan fundamental pada ekosistem
IoT, khususnya terkait penggunaan kata sandi bawaan pabrik (factory default
passwords) dan layanan jaringan yang tidak diamankan.

Neshenko \textit{et al.}~\cite{b8} dalam survei ekstensif mengidentifikasi
bahwa kerentanan pada perangkat IoT dapat dikategorikan ke dalam empat lapisan
utama: lapisan perangkat, lapisan jaringan, lapisan cloud, dan lapisan
kecerdasan buatan. Temuan empiris mereka menunjukkan bahwa sebagian besar
eksploitasi IoT berskala internet menargetkan kerentanan pada lapisan jaringan,
seperti port terbuka yang tidak diperlukan, protokol komunikasi yang tidak
terenkripsi, dan mekanisme autentikasi yang lemah.

Di sisi lain, penelitian Alrawi \textit{et al.}~\cite{b2} menyoroti kesenjangan
kritis antara kompleksitas ancaman keamanan IoT dan kapabilitas teknis pengguna
rumahan. Evaluasi sistematis mereka mengungkapkan bahwa mayoritas pengguna
tidak memiliki pengetahuan, peralatan, maupun kesadaran yang memadai untuk
mengidentifikasi dan memitigasi kerentanan pada perangkat IoT di jaringan rumah
mereka. Temuan ini diperkuat oleh laporan Bitdefender~\cite{b3} yang mencatat
bahwa rata-rata rumah tangga yang terkoneksi memiliki lebih dari 20 perangkat
IoT, namun kurang dari 15\% pengguna yang pernah melakukan pemeriksaan keamanan
terhadap jaringan rumahnya.

Kesenjangan ini menimbulkan urgensi untuk mengembangkan solusi pemindaian
keamanan jaringan yang mudah digunakan (user-friendly) oleh pengguna
non-teknis. Hafeez \textit{et al.}~\cite{b4} mendemonstrasikan bahwa deteksi
ancaman pada tingkat jaringan (network-level detection) merupakan pendekatan
yang paling efektif dan non-intrusive untuk mengidentifikasi aktivitas
berbahaya pada perangkat IoT, tanpa memerlukan modifikasi pada perangkat itu
sendiri. Pendekatan ini menjadi landasan konseptual bagi pengembangan aplikasi
pemindai keamanan jaringan rumah dalam penelitian ini.

Berdasarkan uraian latar belakang di atas, penelitian ini mengusulkan
pengembangan sebuah aplikasi pemindai keamanan jaringan rumah berbasis Python
yang diberi nama HomeGuard. Aplikasi ini dirancang untuk mengidentifikasi
kerentanan perangkat IoT pada jaringan Wi-Fi domestik dengan mengacu pada
kerangka kerja OWASP IoT Top 10~\cite{b9} sebagai standar klasifikasi
kerentanan.

\subsection{Rumusan Masalah}
\begin{enumerate}
    \item Bagaimana merancang dan mengimplementasikan aplikasi pemindai
    keamanan jaringan rumah yang mampu mengidentifikasi kerentanan perangkat
    IoT pada jaringan Wi-Fi domestik?
    \item Apa saja jenis kerentanan yang paling umum ditemukan pada perangkat
    IoT dalam jaringan Wi-Fi domestik berdasarkan kerangka acuan OWASP IoT
    Top 10?
    \item Bagaimana efektivitas aplikasi yang dikembangkan dalam mendeteksi
    kerentanan keamanan pada perangkat IoT di lingkungan jaringan rumah?
\end{enumerate}

\subsection{Tujuan Penelitian}
\begin{enumerate}
    \item Merancang dan membangun aplikasi pemindai keamanan jaringan rumah
    berbasis Python yang mengintegrasikan modul network discovery, port
    scanning, service detection, dan pemetaan kerentanan.
    \item Mengidentifikasi dan mengklasifikasikan jenis-jenis kerentanan yang
    terdapat pada perangkat IoT dalam jaringan Wi-Fi domestik berdasarkan
    standar OWASP IoT Top 10.
    \item Mengevaluasi efektivitas aplikasi yang dikembangkan dalam mendeteksi
    kerentanan keamanan pada perangkat IoT melalui serangkaian pengujian pada
    lingkungan jaringan rumah yang sebenarnya (real-world testing).
\end{enumerate}

\subsection{Manfaat Penelitian}
\textbf{Manfaat Praktis:} Menyediakan sebuah tools pemindai keamanan jaringan
yang open-source, mudah digunakan, dan dapat dioperasikan oleh pengguna rumahan
tanpa memerlukan keahlian teknis mendalam untuk mengaudit keamanan jaringan
Wi-Fi dan perangkat IoT miliknya.

\textbf{Manfaat Akademis:} Berkontribusi pada literatur akademik di bidang
keamanan IoT domestik dengan menyajikan bukti empiris mengenai lanskap
kerentanan perangkat IoT pada jaringan rumah tangga di Indonesia, serta
menawarkan pendekatan sistematis untuk penilaian risiko berbasis OWASP IoT
Top 10.

\subsection{Batasan Masalah}
\begin{enumerate}
    \item Pengujian dilakukan secara eksklusif pada jaringan Wi-Fi rumah milik
    peneliti sendiri, sesuai dengan prinsip etika keamanan siber dan ketentuan
    hukum yang berlaku.
    \item Pemindaian berfokus pada analisis tingkat jaringan (network-level
    analysis), mencakup host discovery, port scanning, dan service detection,
    tanpa melakukan analisis firmware perangkat.
    \item Aplikasi dirancang sebagai non-intrusive scanner yang tidak melakukan
    eksploitasi aktif (active exploitation) terhadap kerentanan yang ditemukan.
    \item Klasifikasi kerentanan mengacu pada kerangka kerja OWASP IoT Top 10
    versi 2018~\cite{b9}.
    \item Antarmuka pengguna dikembangkan berbasis web menggunakan framework
    Streamlit.
\end{enumerate}
